Experiment~1 showed that useful detection information was concentrated at the frontopolar sites and that reduced frontal configurations retained substantially more performance than other anatomical subsets. Under the region-mean comparison in Table~\ref{tab:exp1_subset_summary}, the frontal subset achieved macro-$F_1$ values of 66.92\% on Internal and 66.61\% on Cao2018, compared with 69.68\% and 68.91\%, respectively, for the same frontal electrodes scored within the full-montage run. These differences were significant after Bonferroni correction, with $\Delta F_1=-2.76$ percentage points, $p<0.001$, and $r=-0.80$ on Internal and $\Delta F_1=-2.29$ percentage points, $p=0.026$, and $r=-0.47$ on Cao2018; every anatomical subset likewise differed significantly from its corresponding full-montage electrodes in both corpora. The frontal region nevertheless remained far ahead of every other anatomical region, while its region-mean values were distinct from the full-montage headline oracle values of 88.25\% and 80.68\%. Within the full-montage run, the best individual electrode was Fp2 on Internal ($F_1=86.82\%$) and Fp1 on Cao2018 ($F_1=79.68\%$). When the complete pipeline was instead run on one electrode alone, the best single-electrode result was obtained at Fp1 in both corpora, reaching 83.73\% on Internal and 77.65\% on Cao2018. Performance declined sharply for non-frontal subsets, although the subset analysis did not establish that every frontal channel was necessary. This concentration is compatible with the strong projection of blink activity to the prefrontal scalp and with earlier single-channel and forehead-EEG studies~\citep{chang2016detection,ko2020eyeblink,zhang2023method,tran2021detection,valderrama2018automatic,wang2025sliding,shahbakhti2022simultaneous,aoki2023drowsiness,lopezahumada2023hardware}. The findings therefore motivate further evaluation of reduced frontopolar systems for wearable or in-vehicle use, without establishing that such systems will reproduce the retrospective full-montage operating point in deployment.
